% SPDX-License-Identifier: Apache-2.0
\section{A Foreign Module Inside a Netlist}
\label{sec:x-blackbox}

A module written by hand, held by a lowered unit of units as one of
its children. \code{blinker} is Verilog, in
Listing~\ref{lst:x-blinker}, and again VHDL, in \code{blinker.vhd}:
a count that steps once every \code{PERIOD} cycles while enabled,
put out on eight pins that it drives while enabled and leaves
floating otherwise. It has the three things a netlist has to say to
reach such a module: a parameter, a clock pin, and an inout.

The unit \code{Blinker} is that module twice over. For the
simulation it is a model in Rust, an ordinary unit with a loop. For
the netlist it implements \code{Lower} by hand with \code{foreign},
which names the module, its ports by their own names in the order
\code{run} takes them, the value of its parameter, and which clock
drives its clock pin. \code{Top} holds it beside a lowered register
and joins the two as any unit of units does. The pins are a
\code{Pad}: a port that nothing in the design reads or drives, passed
from \code{Top}'s ports to the module's, so that the generated netlist
preserves an \code{inout} connection across the hierarchy. The simulation
transfers no data over this port.

The netlist instantiates \code{blinker} with its parameter and writes
no module for it. The build simulates that netlist together with the
module's own source, under Verilator with \code{blinker.v} and under
nvc with \code{blinker.vhd}, against the trace the model wrote. So
the model and the module are checked against each other, a cycle at
a time, in both languages.

\begin{figure}[htbp]
\centering
\input{blackbox_timing.tex}
\caption{The foreign module's run: the enable, a cycle later through
the register, the module's tick, and the count it puts out.}
\label{fig:x-blackbox}
\end{figure}

\lstinputlisting[caption={\code{blinker.v}, the module written by
hand.},label={lst:x-blinker}]{lib/examples/blinker.v}

\lstinputlisting[caption={\code{ex\_blackbox.rs}. Run with \code{bazel run //lib/examples:ex\_blackbox}.},label={lst:x-blackbox}]{lib/examples/ex_blackbox.rs}

\lstinputlisting[style=ascii,caption={What \code{ex\_blackbox} printed when the build ran it: the enable and the count every fourth cycle, then the lowering, the top and its register, with the module instantiated and not written.},label={lst:x-blackbox-out}]{out_blackbox.txt}
